\section{The critical surrogate-accuracy threshold}
\label{sec:critical-surrogate-threshold}

The perturbation theorem is most transparent for a general target
scale \(\tau\to0\).  The TPC application later sets
\(\tau=1/J\).

\begin{theorem}[One-sided target transfer]
\label{thm:one-sided-target-transfer}
Let \(a_j,b_j\) be nonzero vectors on the same frames of dimension
at least two, let \(\tau_j>0\) satisfy \(\tau_j\to0\), and fix
constants \(C,K\ge0\).  Suppose
\[
 \|a_j-b_j\|_2\le\delta_j\|b_j\|_2.
\]
\begin{equation}
 \rho(b_j)\le C\tau_j,\qquad
 \delta_j\le K\tau_j
\label{eq:one-sided-transfer-hyp}
\end{equation}
implies, for all sufficiently large \(j\),
\begin{equation}
 \boxed{
\rho(a_j)\le(C+K)\tau_j.}
\label{eq:one-sided-transfer-conc}
\end{equation}
If \(C+K>0\), the asymptotic coefficient \(C+K\) cannot be
uniformly reduced using only \eqref{eq:one-sided-transfer-hyp}.
\end{theorem}

\begin{proof}
Apply \eqref{eq:stronger-additive-bounds}.  For sharpness, take
\[
 r_j=C\tau_j,\qquad \delta_j=K\tau_j
\]
for all sufficiently large \(j\), and choose the upper equality
case from \cref{thm:sharp-envelope}.  Since \(r_j,\delta_j\to0\),
the first branch of \eqref{eq:U-envelope} applies and gives
\[
 \frac{\rho(a_j)}{\tau_j}
 =
 C\sqrt{1-K^2\tau_j^2}
 +K\sqrt{1-C^2\tau_j^2}
 \longrightarrow C+K.
\]
Thus no smaller uniform asymptotic coefficient follows from the two
scalar bounds alone.
\end{proof}

\begin{corollary}[Critical \(J^{-1}\) sufficiency]
\label{cor:J-sufficiency}
If
\[
 \rho(b_J)=O(J^{-1}),\qquad
 \frac{\|a_J-b_J\|_2}{\|b_J\|_2}=O(J^{-1}),
\]
then
\[
 \rho(a_J)=O(J^{-1}).
\]
The same statement holds with \(O(J^{-1})\) replaced throughout by
\(J^{-1}X^{o(1)}\) along a critical TPC scale \(J=J_X\).
\end{corollary}

\begin{proof}
Apply \cref{thm:one-sided-target-transfer} with
\(\tau_J=J^{-1}\).  A sum of two subpower multiples of \(J_X^{-1}\)
is again a subpower multiple of \(J_X^{-1}\).
\end{proof}

\begin{theorem}[Two-sided scale transfer with margin]
\label{thm:two-sided-transfer}
Suppose, for fixed \(0<c\le C\) and \(0\le\vartheta<1\),
\begin{equation}
 \frac{c}{J}\le\rho(b_J)\le\frac{C}{J},
 \qquad
 \delta_J\le\frac{\vartheta c}{J}.
\label{eq:two-sided-margin-hyp}
\end{equation}
Then
\begin{equation}
 \boxed{
 \frac{(1-\vartheta)c}{J}
 \le\rho(a_J)\le
 \frac{C+\vartheta c}{J}.}
\label{eq:two-sided-margin-conc}
\end{equation}
Thus a controlled \(O(J^{-1})\) error preserves
\(\rho\asymp J^{-1}\).  In particular,
\(\delta_J=o(J^{-1})\) preserves every fixed two-sided
\(\rho(b_J)\asymp J^{-1}\) law.
\end{theorem}

\begin{proof}
Use
\[
 \rho(b_J)-\delta_J
 \le\rho(a_J)\le
 \rho(b_J)+\delta_J
\]
from \cref{cor:optimal-additive-stability}.
\end{proof}

\begin{remark}[Upper scale versus two-sided scale]
\label{rem:upper-vs-two-sided}
An error \(K/J\) with arbitrary fixed \(K\) always preserves the
upper order \(O(1/J)\).  To preserve a lower bound as well, the
error constant must stay below the surrogate's lower constant.
Saying only ``\(O(1/J)\)'' suppresses this necessary margin.
\end{remark}

\subsection{Sharp countermodels}

\begin{example}[Why \(o(1)\) is not enough]
\label{ex:o1-not-enough}
Let \(e,f\) be orthonormal in \(\C^2\), let
\(\delta_J\to0\), and put
\begin{equation}
 b_J=f,\qquad
 u_J=\delta_Je+\sqrt{1-\delta_J^2}\,f,\qquad
 a_J=\sqrt{1-\delta_J^2}\,u_J.
\label{eq:tangent-countermodel}
\end{equation}
Then
\begin{equation}
 \frac{\|a_J-b_J\|_2}{\|b_J\|_2}
 =\delta_J,\qquad
 \rho(b_J)=0,\qquad
 \rho(a_J)=\delta_J.
\label{eq:tangent-countermodel-values}
\end{equation}
Choosing \(\delta_J=J^{-1/2}\) gives a surrogate with relative error
\(o(1)\) and perfect zero mode, while
\[
 \rho(a_J)=J^{-1/2}\gg J^{-1}.
\]
Thus structure-free \(o(1)\) relative approximation does not
transfer the critical scale.
\end{example}

\begin{remark}[Actual two-point support]
Take
\[
 e=2^{-1/2}(1,1),\qquad
 f=2^{-1/2}(1,-1).
\]
For every sufficiently small positive \(\delta_J\), both coordinates
of \(a_J\) and \(b_J\) are nonzero.  Hence the countermodel already
lives on one common actual support of size two; it is not a padding
artifact.
\end{remark}

\begin{example}[Sharp order \(J^{-1}\)]
\label{ex:J-sharp-order}
In \cref{ex:o1-not-enough}, choose \(\delta_J=K/J\) for a fixed
\(K>0\).  Then
\[
 \rho(b_J)=0,\qquad
 \rho(a_J)=\frac{K}{J}.
\]
Consequently the perturbation contribution enters at exactly the
\(J^{-1}\) scale.  A uniform theorem that tolerates
\(\delta_J\gg J^{-1}\) cannot conclude
\(\rho(a_J)=O(J^{-1})\).
\end{example}

\begin{example}[A lower signal can be erased]
\label{ex:lower-signal-erased}
Fix \(c>0\), let \(r_J=c/J<1\), and set
\[
 b_J=r_Je+\sqrt{1-r_J^2}\,f,\qquad
 a_J=\sqrt{1-r_J^2}\,f.
\]
Then
\[
 \rho(b_J)=\frac{c}{J},\qquad
 \rho(a_J)=0,\qquad
 \frac{\|a_J-b_J\|_2}{\|b_J\|_2}
 =\frac{c}{J}.
\]
Thus an error with the same leading constant as the lower signal can
destroy two-sided comparability.  The margin in
\cref{thm:two-sided-transfer} is not cosmetic.
\end{example}

\begin{theorem}[Sharp order-of-accuracy barrier]
\label{thm:sharp-accuracy-barrier}
Consider all same-frame pairs \((a_J,b_J)\) in dimension at least
two.  Uniform transfer of
\[
 \rho(b_J)=O(J^{-1})
 \quad\Longrightarrow\quad
 \rho(a_J)=O(J^{-1})
\]
from a relative-error hypothesis
\[
 \|a_J-b_J\|_2\le\delta_J\|b_J\|_2
\]
with a prescribed tolerance \(\delta_J\ge0\) holds if
\(\delta_J=O(J^{-1})\), and fails in general whenever
\(J\delta_J\) is unbounded.  In particular, failure can occur even
when \(\delta_J=o(1)\).
\end{theorem}

\begin{proof}
Sufficiency is \cref{cor:J-sufficiency}.  If
\(J\delta_J\) is unbounded, pass to a subsequence on which it tends
to infinity and put
\[
 \varepsilon_J
 =
 \min\left\{\frac12,\sqrt{\frac{\delta_J}{J}}\right\}.
\]
Eventually \(\delta_J>J^{-1}\), so
\(\varepsilon_J\le\delta_J\), while
\[
 J\varepsilon_J
 =
 \min\left\{\frac J2,\sqrt{J\delta_J}\right\}
 \longrightarrow\infty.
\]
Use the tangent construction of \cref{ex:o1-not-enough} with actual
relative error \(\varepsilon_J\) on that subsequence.  Then
\(\rho(b_J)=0\), whereas
\[
 J\rho(a_J)=J\varepsilon_J
\]
is unbounded.  The remaining indices may be filled by any exact
pair \(a_J=b_J\perp e\).
\end{proof}

\begin{remark}[What the barrier does not say]
The theorem is a uniform worst-case statement.  A special surrogate
may transfer with weaker norm accuracy if one separately proves
that its error is almost orthogonal to the constant mode, or proves
a signed estimate for \(Z(a-b)\) stronger than Cauchy--Schwarz.
Such information is additional structure, not a consequence of
\(\delta=o(1)\).
\end{remark}
